% !TeX root = ./rpz.tex
\chapter{{\fontsize{14pt}{21pt}\selectfont\MakeUppercase{Реализация}}}
\label{cha:design}

\section{Модель данных}

На рисунке~\ref{fig:database_model} представлена логическая модель данных разрабатываемого веб-приложения.

\begin{figure}[H]
  \centering
  \includegraphics[width=0.9\textwidth]{database_model.png}
  \caption{Модель данных}
  \label{fig:database_model}
\end{figure}

Структура базы данных включает следующие функциональные группы сущностей:

\begin{compactlist}
  \item хранение информации о пользователях;
  \item хранение задач с географической привязкой;
  \item хранение маршрутов;
  \item хранение результатов вычислений.
\end{compactlist}

Сущности, связанные с пользователями. Сущность users предназначена для хранения учетных записей пользователей.
Первичный ключ id используется для установления связей со всеми зависимыми таблицами.
Поле email выполняет роль логина и имеет ограничение уникальности, что исключает дублирование учетных записей.
Пароль хранится в виде хэшированного значения, что обеспечивает безопасность аутентификации.

Сущность refresh\_tokens используется для поддержки механизма продления авторизации.
Каждый токен связан с конкретным пользователем через внешний ключ user\_id.
Хранение токена в виде хэша предотвращает его компрометацию в случае несанкционированного доступа к базе данных.

Сущность задач. Сущность tasks отражает основную предметную сущность системы — задачу, подлежащую выполнению.

Каждая запись содержит:

\begin{compactlist}
  \item связь с владельцем задачи (user\_id);
  \item текстовое описание адреса (address\_text);
  \item географические координаты (latitude, longitude), допускающие NULL для задач без адреса;
  \item длительность выполнения (duration\_min), допускающую NULL для мгновенных задач;
  \item раздельные дату и время начала (window\_start\_date, window\_start\_time) и окончания (window\_end\_date, window\_end\_time) временного окна, что позволяет задавать ограничения как по дате, так и по времени;
  \item индекс сортировки (sort\_index) для отображения порядка в пользовательском интерфейсе;
  \item признак выполнения (is\_completed) для отметки завершенных задач.
\end{compactlist}

Сущности маршрутов. Сущность routes описывает факт построения маршрута для конкретного пользователя.
На каждого пользователя хранится не более одного маршрута, что обеспечивается ограничением уникальности UNIQUE(user\_id): при повторном запуске оптимизации предыдущая запись и связанные с ней остановки удаляются каскадно.
Запись содержит идентификатор использованного алгоритма (algorithm), агрегированные показатели маршрута (суммарную длину total\_distance\_m, суммарное время перемещения total\_travel\_sec, суммарное время выполнения задач total\_service\_sec и суммарное ожидание total\_wait\_sec) и момент вычисления computed\_at.
Размещение агрегированных метрик непосредственно в строке маршрута упрощает чтение результатов оптимизации одним запросом и устраняет необходимость отдельной таблицы статистики.

Структура маршрута детализируется в сущности route\_stops,
где каждая запись соответствует отдельной точке маршрута.
Поле position определяет порядок следования,
а временные поля (arrive\_time, service\_start\_time, service\_end\_time)
позволяют зафиксировать результат работы алгоритма с учетом временных окон.
Каскадное удаление по внешним ключам route\_id и task\_id гарантирует согласованность данных при пересчете маршрута или удалении задачи.

Геометрия маршрута. Геометрическое представление маршрута (полилинии участков между остановками) на сервере не хранится: оно строится клиентской частью с использованием Yandex Maps JS API на основе координат остановок и выбранного режима передвижения.
Такой подход исключает необходимость отдельной таблицы для геометрии и обеспечивает актуальность визуализации при изменении транспортного режима без обращения к серверу.

\section{Архитектура системы}

Состав и взаимодействие компонентов. Система построена по клиент-серверной
модели и состоит из трех развертываемых компонентов: клиентского
одностраничного приложения (SPA) на React, серверного REST API на Go и
реляционной базы данных PostgreSQL. Расчет матрицы расстояний и отрисовка
карты выполняются через внешний сервис Yandex Maps JS API, а резервный
расчет матрицы~--- через сервис OSRM. Общая структура контейнеров и связи
между ними в нотации C4 (уровень контейнеров) приведены
на рисунке~\ref{fig:c4_container}.

\begin{figure}[H]
  \centering
  \includegraphics[width=\textwidth]{c4_container.png}
  \caption{Диаграмма контейнеров системы в нотации C4}
  \label{fig:c4_container}
\end{figure}

Браузер пользователя загружает SPA, которое обращается к серверу по
протоколу HTTPS запросами REST; авторизация выполняется по JWT токену.
Сервер обращается к PostgreSQL и инкапсулирует бизнес-логику: аутентификацию,
управление задачами и оптимизацию маршрута. Матрица расстояний строится
преимущественно на стороне клиента средствами Yandex Maps JS API,
поскольку бесплатная маршрутизация Яндекса доступна только из JavaScript;
если клиент не передал готовую матрицу~--- например, при недоступности Яндекса
после повторных попыток или при обращении к API в обход веб-клиента,~--- сервер
рассчитывает ее самостоятельно через OSRM.
Такое разделение ответственности минимизирует серверные обращения к
платным сервисам и переносит интерактивную работу с картой на клиента.

Общие принципы. Серверная часть построена по принципам чистой архитектуры, предложенной Р.~Мартином~\cite{martin2017clean}.
Зависимости направлены от внешних слоев (HTTP-обработчики, репозитории)
к внутренним (доменные сущности и варианты использования), но не наоборот.
Такое разделение делает бизнес-логику независимой от конкретного
HTTP-фреймворка, СУБД и сторонних сервисов.

Управление задачами. Ядро прикладной логики составляет работа со списком
задач пользователя~--- именно сформированный список служит входными данными
для алгоритма оптимизации маршрута. Доменная сущность задачи содержит
название, текстовый адрес и географические координаты, длительность
выполнения, временное окно (раздельные дата и время начала и окончания)
и порядковый индекс в списке. Варианты использования реализуют полный набор
операций над списком: создание и редактирование задачи, удаление, получение
перечня, пакетное создание, импорт из файла и изменение порядка следования.
Доступ к данным инкапсулирован в репозитории поверх PostgreSQL в соответствии
с принятым разделением на слои: HTTP-обработчики принимают и валидируют
запросы, варианты использования содержат бизнес-правила, репозиторий выполняет
операции с базой данных. Каждая операция выполняется в контексте
аутентифицированного пользователя, что изолирует списки задач разных
пользователей друг от друга.

Кэширование внешних запросов. Обращения к внешним сервисам маршрутизации
кэшируются, чтобы при повторной оптимизации одного и того же набора задач
не выполнять одинаковые запросы заново. Кэш разделен на два независимых
уровня~--- клиентский и серверный, что обусловлено способом доступа
к каждому из сервисов маршрутизации.

Клиентский кэш. Матрица расстояний строится в браузере средствами
Yandex Maps JS API, поскольку бесплатная маршрутизация Яндекса доступна
только из JavaScript; серверный аналог предоставляется
лишь по платной коммерческой лицензии и в проекте не используется. По этой
причине кэш рассчитанных ребер также размещен на стороне клиента~--- в
оперативной памяти в пределах пользовательской сессии. В качестве хранилища
применяется LRU-кэш с ограничением размера и временем жизни записи 24~часа.
Стратегия LRU выбрана как одна из наиболее эффективных и наименее ресурсоемких политик
вытеснения для кэшей с ограниченным объемом~\cite{kuduhov2014lru}.
Перед обращением к Yandex API матрица заполняется из
кэша, и во внешний сервис отправляются только отсутствующие пары точек, после
чего рассчитанные значения сохраняются обратно. Ключом записи служит режим
передвижения и упорядоченная пара координат, округленных до пятого знака после
запятой; сравнение выполняется по строковому
представлению ключа, что исключает влияние погрешности вычислений с плавающей
точкой и незначительных колебаний координат, возвращаемых геокодером.

Серверный кэш. Результаты обращений к серверному сервису маршрутизации OSRM
(резервный путь, когда клиент не передал готовую матрицу) кэшируются на уровне
приложения в оперативной памяти через обертку CachedProvider над интерфейсом
distance.Provider. В качестве хранилища используется библиотека
hashicorp/golang-lru/v2~\cite{golang_lru} (пакет expirable)~--- потокобезопасный
LRU-кэш с ограничением размера и временем жизни записи (TTL). Ключом записи
служит пара географических координат, округленных до шестого знака после запятой,
что делает кэш инвариантным к порядку запросов одних и тех же адресов.

Оба уровня кэша построены на единой политике вытеснения LRU и снижают количество
обращений к внешним сервисам, уменьшают задержки повторной оптимизации и повышают
устойчивость системы при ограничениях со стороны API-провайдеров.

\section{Алгоритм оптимизации маршрута}

Описание алгоритма. Для вычисления оптимального порядка посещения задач реализован алгоритм
«ближайшего соседа с временными окнами» (Nearest Neighbour Heuristic with
Time Windows, NNH-TW).
Базовый алгоритм ближайшего соседа~--- одна из наиболее изучаемых
конструктивных эвристик для задачи коммивояжера, традиционно
используемая в качестве эталона при сравнении более сложных
методов~\cite{garey1982intractable}.

Алгоритм является жадной конструктивной эвристикой с одношаговым просмотром
вперед: на каждом шаге из допустимых кандидатов выбирается
узел с наименьшим временем завершения обслуживания, при этом предпочтение
отдается <<безопасным>> кандидатам --- тем, посещение которых не делает
недостижимым ни один другой узел с временным окном. Внутри каждого класса
приоритет получает кандидат с более ранним
дедлайном, что снижает вероятность пропуска срочных задач.
Алгоритм также учитывает дополнительные ограничения: фиксацию начальной
и конечной точки маршрута, а также попарные ограничения порядка посещения.

Псевдокод алгоритма приведен в листинге~\ref{listing:nnh_pseudocode}.

\begin{codelisting}{Псевдокод алгоритма NNH-TW с просмотром вперед}{listing:nnh_pseudocode}
  \begin{minted}{text}
Вход: граф G(V, E), startTimeUnix, constraints{startIdx, endIdx, precedences}
Выход: порядок обхода order[], тайминги[], агрегированная статистика

1.  prereqs <- buildPrereqs(precedences)
2.  endIdx <- constraints.endIdx (или -1)
3.  Если constraints.startIdx задан:
        cur <- constraints.startIdx
    Иначе:
        cur <- узел с наиболее ранним окном среди допустимых
4.  visited[cur] <- true;  order <- [cur]
5.  currentTimeSec <- startTimeUnix + duration[cur] * 60

6.  Пока |order| < |V|:
    // Если остался только закрепленный конечный узел — ставим его
    a. Если endIdx >= 0 И все прочие посещены:
         Добавить endIdx в маршрут; выйти из цикла

    // Проход 1: допустимый узел с минимальным completionTime + look-ahead
    b. bestComp <- +inf;  bestSafe <- false;  bestDeadline <- +inf;  next <- -1
    c. Для каждого i not in visited, i != endIdx:
         Если не выполнены предшественники i: пропустить
         arrivalSec <- currentTimeSec + travelSec[cur][i]
         Если НЕ feasible(i, arrivalSec): пропустить
         comp <- completionTime(i, arrivalSec)
         safe <- НЕ causesWindowMiss(i, comp, G, visited, endIdx)
         Если (safe > bestSafe) ИЛИ (safe = bestSafe И comp < bestComp)
              ИЛИ (safe = bestSafe И comp = bestComp И deadline[i] < bestDeadline):
             next <- i;  bestComp <- comp;  bestSafe <- safe

    // Проход 2: резервный — минимальный completionTime без учета окна
    d. Если next = -1:
         Для каждого i not in visited, i != endIdx:
             Если не выполнены предшественники i: пропустить
             comp <- completionTime(i, arrivalSec)
             Если comp < bestComp: next <- i

    e. waitSec <- max(0, windowStart[next] - arrivalSec)
    f. currentTimeSec <- arrivalSec + waitSec + duration[next] * 60
    g. visited[next] <- true;  order.append(next);  cur <- next

7.  Вернуть order, тайминги и статистику

feasible(i, arrivalSec):
    Если windowEnd[i] < 0: вернуть true
    Вернуть arrivalSec <= windowEnd[i] - duration[i] * 60

causesWindowMiss(cand, afterCandSec, G, visited, endIdx):
    Для каждого j not in visited, j != cand, j != endIdx:
        Если windowEnd[j] < 0: пропустить
        Если НЕ feasible(j, afterCandSec + travelSec[cand][j]):
            Вернуть true
    Вернуть false
  \end{minted}
\end{codelisting}

Блок-схема алгоритма разбита на два рисунка. На
рисунке~\ref{fig:nnh_tw_main} показан основной цикл построения маршрута.
Финишная задача и обязательный порядок выполнения (когда одна задача
должна предшествовать другой) задаются пользователем как ограничения.
Каждая задача может иметь временное окно~--- интервал, в который ее
допустимо выполнить. Стартовую задачу пользователь также может указать явно, иначе она
выбирается эвристически~--- задача без предшественников с наиболее ранним
временным окном. Финишная задача, если она указана, резервируется и
обслуживается последней: она исключается из обычного выбора и
добавляется в конец, когда все прочие задачи уже размещены. На рисунке~\ref{fig:nnh_tw_select}
приведена подпрограмма выбора очередной задачи. Проход~1 с просмотром
вперед берет задачу, до которой успеваем в ее временное окно и все
обязательные предшественники которой уже выполнены, отдавая предпочтение
той, после которой остальные задачи все еще успевают в свои окна и
которая освобождает исполнителя раньше. Резервный проход~2 включается,
если в проходе~1 подходящих задач не нашлось, и берет ближайшую по
времени завершения задачу без учета временных окон, но также соблюдая
обязательный порядок выполнения.

\begin{figure}[H]
  \centering
  \includegraphics[width=0.92\textwidth]{nnh_tw_main.png}
  \caption{Блок-схема алгоритма NNH-TW: основной цикл построения маршрута}
  \label{fig:nnh_tw_main}
\end{figure}

\begin{figure}[H]
  \centering
  \includegraphics[width=0.64\textwidth]{nnh_tw_select.png}
  \caption{Блок-схема подпрограммы выбора следующей задачи}
  \label{fig:nnh_tw_select}
\end{figure}

Трассировка на числовом примере. Ниже приведен пример с тремя задачами для иллюстрации работы алгоритма.
Исходные данные приведены в таблице~\ref{tab:trace_input},
матрица времен перемещения — в таблице~\ref{tab:trace_matrix}.

\begin{table}[H]
  \caption{Исходные данные задач для трассировки}
  \label{tab:trace_input}
  \begin{tabularx}{\linewidth}{|C|L|C|C|}
    \hline
    Узел & \multicolumn{1}{c|}{Адрес} & Временное окно & Длительность \\ \hline
    $v_0$                   & Офис (начало)            & 08:00--09:00                      & 30~мин                                     \\ \hline
    $v_1$                   & Точка А                  & 10:00--12:00                      & 20~мин                                     \\ \hline
    $v_2$                   & Точка Б                  & 09:00--10:00                      & 15~мин                                     \\ \hline
  \end{tabularx}
\end{table}

\begin{table}[H]
  \caption{Матрица времен перемещения (минуты)}
  \label{tab:trace_matrix}
  \begin{tabularx}{\linewidth}{|C|C|C|C|}
    \hline
    Из / В & $v_0$ & $v_1$ & $v_2$ \\ \hline
    $v_0$           & ---   & 40    & 15    \\ \hline
    $v_1$           & 40    & ---   & 20    \\ \hline
    $v_2$           & 15    & 20    & ---   \\ \hline
  \end{tabularx}
\end{table}

Трассировка алгоритма при начале работы в 08:00 (для наглядности
время приведено в минутах от полуночи; в реализации используются
секунды~Unix):

Выбор стартового узла. Узел $v_0$ имеет наиболее раннее
окно (08:00), поэтому именно он выбирается стартовым:
$\mathrm{cur} = 0$, $\mathrm{currentTime} = 480 + 30 = 510$
(08:30 плюс 30~мин выполнения, что соответствует 09:30 или 510~мин от полуночи).

Шаг 1. Проход~1. Для кандидата $v_1$ прибытие составляет
$510 + 40 = 550$ (09:10). Окно $v_1$ задано как 600--720 (10:00--12:00),
условие $550 \leq 720 - 20 = 700$ выполнено, время в пути $t = 40$~мин.
Для кандидата $v_2$ прибытие составляет $510 + 15 = 525$ (08:45),
окно 540--600 (09:00--10:00), условие $525 \leq 600 - 15 = 585$
выполнено, $t = 15$~мин. Поскольку $v_2$ ближе ($15 < 40$),
выбирается $\mathrm{next} = 2$.
Ожидание открытия окна составляет $\max(0,\; 540 - 525) = 15$~мин,
что дает $\mathrm{currentTime} = 525 + 15 + 15 = 555$ (09:15).

Шаг 2. Проход~1. Для оставшегося кандидата $v_1$ прибытие
составляет $555 + 20 = 575$ (09:35), условие $575 \leq 700$ выполнено,
$\mathrm{next} = 1$. Ожидание равно $\max(0,\; 600 - 575) = 25$~мин,
после чего $\mathrm{currentTime} = 575 + 25 + 20 = 620$ (10:20).

Итоговый порядок посещения узлов: $v_0 \to v_2 \to v_1$. Суммарное
ожидание открытия временных окон по маршруту составляет
$15 + 25 = 40$~мин. Полученная последовательность соответствует
очередности, при которой ни одно из заданных временных окон не нарушается,
а суммарное время в пути и время ожидания минимизированы.

Оценка вычислительной сложности. Алгоритм выполняет $n$ итераций основного цикла; на каждой итерации
для каждого из $n$ кандидатов выполняется одношаговый просмотр вперед
(causesWindowMiss) по $n$ оставшимся узлам, что дает
сложность $O(n^3)$.

При $n = 30$ задачах алгоритм выполняет порядка $30^3 = 27\,000$ элементарных
операций. Результаты бенчмаркинга, приведенные в таблице~\ref{tab:algo_perf},
показывают, что при $n \leq 30$ время работы алгоритма составляет единицы
микросекунд, что на несколько порядков меньше задержки, вносимой получением
матрицы расстояний от внешнего API.

Дополнительные ограничения маршрута. Помимо временных окон, реализована поддержка трех типов дополнительных ограничений,
задаваемых пользователем при запросе оптимизации.

Фиксированная начальная точка.
Если пользователь указывает идентификатор начальной задачи, алгоритм начинает обход
с соответствующего узла вместо автоматически выбранного по признаку наиболее раннего
временного окна.

Фиксированная конечная точка.
Если задана конечная задача, соответствующий узел исключается из стандартного процесса
выбора следующей точки.
После того как все остальные узлы посещены, конечный узел добавляется в маршрут последним.

Ограничения порядка.
Пользователь вправе указать пары задач вида «A до B»:
задача $A$ должна быть выполнена строго раньше задачи $B$.
При выборе следующего узла на каждом шаге алгоритм пропускает кандидатов,
для которых не все предшественники уже посещены.
Условие проверяется в обоих проходах (Pass~1 и Pass~2), что гарантирует
соблюдение ограничения независимо от наличия допустимых временных окон.

Ограничения передаются в алгоритм через структуру Constraints пакета
optimizer.
Преобразование идентификаторов задач в индексы графа выполняется на уровне
бизнес-логики (story.Optimize) и не затрагивает реализацию алгоритма,
что соответствует принципу единственной ответственности.

\section{Интерфейс}

На рисунках~\ref{fig:maket_screen_login}, \ref{fig:app_screen_main}
и~\ref{fig:app_screen_optimized} приведены экранные формы реализованного приложения.
Рисунок~\ref{fig:maket_screen_login} соответствует странице авторизации,
рисунок~\ref{fig:app_screen_main}~--- главной странице до запуска оптимизации,
рисунок~\ref{fig:app_screen_optimized}~--- ей же после построения маршрута.

\begin{figure}[H]
  \centering
  \includegraphics[width=0.85\textwidth]{maket_screen_login.png}
  \caption{Страница авторизации приложения}
  \label{fig:maket_screen_login}
\end{figure}
\vspace{5pt}

\begin{figure}[H]
  \centering
  \includegraphics[width=\textwidth]{app_screen_main.png}
  \caption{Главная страница приложения до запуска оптимизации}
  \label{fig:app_screen_main}
\end{figure}
\vspace{5pt}

\begin{figure}[H]
  \centering
  \includegraphics[width=\textwidth]{app_screen_optimized.png}
  \caption{Главная страница после построения маршрута}
  \label{fig:app_screen_optimized}
\end{figure}

Структура и организация. Главная страница разделена на две функциональные зоны (рисунок~\ref{fig:app_screen_main}):
слева~--- список задач, справа~--- карта с визуализацией маршрута.
Размещение в пределах одного экрана позволяет немедленно оценивать влияние
изменений в списке задач на итоговый маршрут без переключения между
представлениями, что важно для пользователей без специальной логистической
подготовки. Список задач отображается в виде отдельных карточек
с фиксированным набором полей (номер, название, адрес, длительность,
временное окно). Карта реализована на платформе Яндекс Карты, что
обеспечивает связь между порядком задач и их пространственным расположением.

После запуска оптимизации (рисунок~\ref{fig:app_screen_optimized}) над
списком появляется панель сводки маршрута с суммарным расстоянием,
временем в пути и общим временем выполнения; на карте~--- линия маршрута
с пронумерованными метками и переключатель способа передвижения; карточки
перестраиваются в порядке посещения, рассчитанном алгоритмом.

Добавление задачи выполняется через модальное окно с логически
сгруппированными полями: название, адрес, длительность, временное окно
(дата и время начала, дата и время окончания) и роль задачи в маршруте
(рисунок~\ref{fig:app_ui_add_task}). Обязательно только название; адрес,
длительность и окно необязательны, что позволяет вводить как полноценные
точки маршрута, так и задачи без географической привязки; минимизация числа
обязательных полей снижает порог входа для частных пользователей без
специальной подготовки. Поле адреса
поддерживает автодополнение через сервис геокодирования Яндекса: по мере
ввода под полем отображается список вариантов, выбор которого подставляет
координаты задачи. Кнопки <<Начальная точка>> и <<Конечная точка>>
позволяют сразу при создании закрепить задачу как старт или финиш маршрута.

\begin{figure}[H]
  \centering
  \includegraphics[width=0.52\textwidth]{app_ui_add_task.png}
  \caption{Модальное окно добавления задачи с автодополнением адреса}
  \label{fig:app_ui_add_task}
\end{figure}
\vspace{5pt}

Навигация организована
через верхнюю панель и включает две вкладки~--- <<Главная>> и
<<Инструкция>>; в правой части шапки размещены идентификатор активной
учетной записи и кнопка выхода из сессии.

Импорт и экспорт задач. Для ускорения ввода большого количества задач реализован импорт из файлов
CSV, JSON и Excel~(XLSX). Диалог импорта содержит описание ожидаемого формата
(столбцы названия, адреса, длительности, даты и границ временного окна) и
кнопки загрузки шаблона, а после выбора файла строит предпросмотр
распознанных строк (рисунок~\ref{fig:app_ui_import}). Каждая строка проверяется:
подсчитывается число корректных и ошибочных записей, а строки с нарушениями
(пустое название, нечисловая длительность, неверный формат времени, конец
окна раньше начала) подсвечиваются красным с пояснением причины. Строки с
ошибками пропускаются при сохранении, корректные~--- импортируются и
геокодируются автоматически, что исключает попадание некорректных данных в
список задач.

\begin{figure}[H]
  \centering
  \includegraphics[width=\textwidth]{app_ui_import.png}
  \caption{Предпросмотр импортируемых задач с проверкой строк на ошибки}
  \label{fig:app_ui_import}
\end{figure}

Экспорт задач доступен в двух форматах: CSV и Excel. Экспортируются все текущие задачи
пользователя с полями: название, адрес, длительность, дата и время временного окна.

Дополнительные элементы управления. В интерфейсе реализована возможность изменения порядка задач методом
перетаскивания: даже после автоматической оптимизации пользователь может
скорректировать последовательность вручную без повторного запуска алгоритма.
Помимо этого, реализованы перечисленные ниже элементы управления.

Выбор режима передвижения. Над картой размещен переключатель способа
перемещения~--- пешком, общественный транспорт или автомобиль
(рисунок~\ref{fig:app_ui_transport}). Выбранный режим определяет матрицу
расстояний, запрашиваемую у Яндекс Карт, и влияет на визуализацию: для
автомобильного и пешего маршрута строится единая линия с суммарным временем
и расстоянием, а для общественного транспорта~--- цепочка
сегментов с указанием номеров маршрутов и остановок.

\begin{figure}[H]
  \centering
  \includegraphics[width=0.62\textwidth]{app_ui_transport.png}
  \caption{Переключатель режима передвижения и визуализация маршрута на карте}
  \label{fig:app_ui_transport}
\end{figure}
\vspace{5pt}

Фиксация начальной и конечной точек. Через контекстное меню карточки задачи
(рисунок~\ref{fig:app_ui_context_menu}) пользователь назначает точку старта
или финиша маршрута. Закрепленные задачи помечаются в списке метками
<<Старт>> и <<Финиш>> (рисунок~\ref{fig:app_ui_start_end}) и передаются
алгоритму как фиксированные концевые узлы.

\begin{figure}[H]
  \centering
  \includegraphics[width=0.85\textwidth]{app_ui_context_menu.png}
  \caption{Контекстное меню карточки задачи: назначение старта и финиша,
           клонирование, редактирование и удаление}
  \label{fig:app_ui_context_menu}
\end{figure}
\vspace{5pt}

\begin{figure}[H]
  \centering
  \includegraphics[width=0.7\textwidth]{app_ui_start_end.png}
  \caption{Список задач с отмеченными точками старта и финиша и временными окнами}
  \label{fig:app_ui_start_end}
\end{figure}
\vspace{5pt}

Ограничения порядка. Пользователь задает пары <<задача A до задачи B>>
(рисунок~\ref{fig:app_ui_precedence}), которые алгоритм обязан соблюдать при
построении маршрута; число таких пар не ограничено.

\begin{figure}[H]
  \centering
  
\includegraphics[width=\textwidth]{app_ui_precedence.png}
  \caption{Задание ограничений порядка выполнения задач}
  \label{fig:app_ui_precedence}
\end{figure}
\vspace{5pt}

Отметка выполнения. Задачу можно пометить как выполненную; такая карточка
отображается приглушенным цветом с зачеркнутым названием
(рисунок~\ref{fig:app_ui_complete}) и исключается из последующей оптимизации.

\begin{figure}[H]
  \centering
  \includegraphics[width=0.7\textwidth]{app_ui_complete.png}
  \caption{Задача, отмеченная как выполненная}
  \label{fig:app_ui_complete}
\end{figure}
\vspace{5pt}

Клонирование задач. Команда <<Клонировать>> контекстного меню создает копию
задачи с суффиксом <<(копия)>> и теми же адресом, длительностью и временным
окном (рисунок~\ref{fig:app_ui_clone}), что ускоряет ввод однотипных записей.

\begin{figure}[H]
  \centering
  
\includegraphics[width=0.7\textwidth]{app_ui_clone.png}
  \caption{Результат клонирования задачи}
  \label{fig:app_ui_clone}
\end{figure}
\vspace{5pt}

Обнаружение конфликтов. Интерфейс автоматически выявляет задачи с
некорректными или взаимно несовместимыми временными окнами и работает в двух
режимах~--- до и после оптимизации.

До оптимизации проверяются сами окна. Помечаются задачи с инвертированным
окном (время окончания не позже времени начала), задачи, длительность
обслуживания которых превышает ширину окна, а также группы задач, чьи окна
перекрываются настолько, что суммарное время обслуживания не умещается в
общий интервал. Конфликтующие карточки выделяются красной рамкой, красной
меткой окна и значком предупреждения
(рисунок~\ref{fig:app_ui_conflict_before}), а кнопка запуска оптимизации
блокируется с подсказкой об устранении конфликтов. Тем самым исключается
запуск алгоритма для заведомо неразрешимой постановки.

\begin{figure}[H]
  \centering
  \includegraphics[width=0.7\textwidth]{app_ui_conflict_before.png}
  \caption{Обнаружение конфликтов временных окон до запуска оптимизации}
  \label{fig:app_ui_conflict_before}
\end{figure}
\vspace{5pt}

После оптимизации проверяется фактическое расписание. Для каждой остановки
расчетное время прибытия сопоставляется с крайним сроком окна с учетом
длительности обслуживания. Если задача не успевает быть обслужена в пределах
своего окна, над списком остановок выводится предупреждение о числе таких
задач, а соответствующая остановка подсвечивается красным с отметкой
<<Не укладывается в окно>> (рисунок~\ref{fig:app_ui_conflict_after}). Это
позволяет увидеть, что при заданном наборе ограничений отдельные окна
физически недостижимы, и скорректировать исходные данные~--- длительность,
границы окна или состав маршрута.

\begin{figure}[H]
  \centering
  \includegraphics[width=0.92\textwidth]{app_ui_conflict_after.png}
  \caption{Обнаружение фактического опоздания после построения маршрута}
  \label{fig:app_ui_conflict_after}
\end{figure}

Адаптивная верстка. Интерфейс сохраняет работоспособность
на устройствах различных классов: настольные мониторы, планшеты и смартфоны;
поддержка смартфонов существенна для курьеров и мастеров на выезде, работающих
вне стационарного рабочего места.
Адаптация выполнена средствами фреймворка
Tailwind~CSS~\cite{tailwind_responsive} с применением контрольных точек ширины.

На широких экранах (рисунок~\ref{fig:app_screen_main}) список задач и карта
размещаются рядом в две колонки, что обеспечивает их одновременный обзор.
При сужении области просмотра до планшета (рисунок~\ref{fig:app_screen_tablet})
двухколоночная компоновка перестраивается в одноколоночную: карта сдвигается
под список задач, а кнопки верхней панели инструментов располагаются
в две строки. На экранах смартфонов (рисунок~\ref{fig:app_screen_mobile})
текстовые подписи кнопок сворачиваются до иконок, а длинные названия задач
усекаются с многоточием для сохранения целостности карточек.

\begin{figure}[H]
  \centering
  \includegraphics[width=0.7\textwidth]{app_screen_tablet.png}
  \caption{Главная страница приложения на планшете}
  \label{fig:app_screen_tablet}
\end{figure}
\vspace{5pt}

\begin{figure}[H]
  \centering
  \includegraphics[width=0.45\textwidth]{app_screen_mobile.png}
  \caption{Главная страница приложения на смартфоне после построения маршрута}
  \label{fig:app_screen_mobile}
\end{figure}

\section{Техническая реализация}

Состав и объем программного кода. Приложение включает два независимых программных компонента: серверную часть
на Go и клиентскую часть на TypeScript с библиотекой React. Серверная часть
содержит около 54 файлов Go, организованных по пакетам в соответствии
с принципами чистой архитектуры (доменные сущности и варианты использования,
HTTP-обработчики, репозитории PostgreSQL, провайдер матрицы расстояний)
и пять SQL-миграций. Клиентская часть включает около 100 файлов
TypeScript/TSX, из которых около 40 реализуют прикладную логику и
компоненты страниц, а остальные~--- примитивы библиотеки пользовательского
интерфейса на основе shadcn/ui~\cite{shadcn} и Radix~UI~\cite{radix_ui}.

Наиболее значимые фрагменты кода приведены в приложениях: реализация
алгоритма оптимизации маршрута~--- в приложении~\ref{app:lst_optimizer},
HTTP-обработчики маршрутов~--- в приложении~\ref{app:lst_handlers_routes},
клиент API с логикой обновления токена~--- в приложении~\ref{app:lst_api_client},
утилита построения матрицы расстояний~--- в приложении~\ref{app:lst_route_optimizer}.

REST-API управления задачами. Операции планирования задач доступны клиенту
через набор HTTP-эндпоинтов с аутентификацией по JWT токену; тело запросов и
ответов передается в формате JSON:

\begin{compactlist}
  \item GET /api/tasks~--- получение списка задач пользователя;
  \item POST /api/tasks~--- создание задачи;
  \item POST /api/tasks/batch~--- пакетное создание нескольких задач за один запрос;
  \item POST /api/tasks/import~--- импорт задач из файла формата CSV или XLSX;
  \item PATCH /api/tasks/order~--- изменение порядка следования задач в списке;
  \item PATCH /api/tasks/\{id\}~--- частичное обновление полей задачи;
  \item DELETE /api/tasks/\{id\}~--- удаление задачи.
\end{compactlist}

Требования к техническим средствам. Для развертывания и эксплуатации приложения предъявляются следующие
требования к техническим средствам.

Для серверной части:

\begin{compactlist}
  \item процессор: одноядерный с тактовой частотой не ниже 1~ГГц
        (рекомендуется двухъядерный и более);
  \item оперативная память: не менее 512~МБ (рекомендуется 1~ГБ и более);
  \item дисковое пространство: не менее 1~ГБ для хранения образов Docker
        и данных PostgreSQL;
  \item сетевой интерфейс: соединение с интернетом для обращения к внешним
        сервисам (Яндекс Карты, OSRM).
\end{compactlist}

Для клиентской части (рабочее место пользователя):

\begin{compactlist}
  \item устройство с современным веб-браузером и поддержкой JavaScript;
  \item соединение с интернетом для загрузки Яндекс Карт и обращения к API.
\end{compactlist}

Требования к программному обеспечению. Для развертывания серверной части необходимы:

\begin{compactlist}
  \item операционная система Linux~\cite{linux_kernel}, macOS~\cite{macos}
        или Windows~\cite{windows} с поддержкой контейнеризации;
  \item Docker Engine версии 24.0 и выше;
  \item Docker Compose версии 2.0 и выше.
\end{compactlist}

Для сборки и запуска клиентской части в режиме разработки необходимы
Node.js~\cite{nodejs} версии 18 и выше и менеджер пакетов
pnpm~\cite{pnpm} версии 8 и выше (либо npm версии 9 и выше). На стороне пользователя поддерживаются
браузеры Google Chrome версии 90 и выше, Mozilla Firefox версии 88
и выше, Apple Safari версии 14 и выше, Microsoft Edge версии 90 и выше.

\section{Выводы}

Спроектирована модель данных, поддерживающая хранение задач, маршрутов,
временных окон и пользовательских ограничений.

Серверная часть построена по принципам чистой архитектуры с разделением
слоев домена, вариантов использования и инфраструктуры; реализован
алгоритм NNH-TW с одношаговым просмотром вперед и поддержкой
фиксированных начальной и конечной точек маршрута, а также попарных
ограничений порядка посещения.

Разработан клиентский интерфейс с интерактивной картой и средствами
ручного управления маршрутом; обеспечена контейнеризация приложения
на основе Docker Compose для воспроизводимого развертывания.
